\title{Large Language Models Can Automatically Engineer Features for Few-Shot Tabular Learning}

\begin{abstract}
Large Language Models (LLMs) have demonstrated exceptional capability in addressing complex, previously unseen reasoning tasks, making them highly promising tools for tabular learning—a cornerstone of numerous practical applications. In this work, we introduce a new in-context learning paradigm, \textsf{FeatLLM}, which leverages LLMs as feature constructors to transform input datasets, thereby facilitating superior tabular prediction. These newly created features are subsequently utilized by a straightforward downstream machine learning model, such as linear regression, enabling highly effective few-shot learning. Importantly, \textsf{FeatLLM} requires only this simple predictive model, together with the engineered features, during inference. Unlike prior LLM-centric approaches, \textsf{FeatLLM} does not necessitate querying the LLM for every test sample during deployment. Furthermore, it operates with only API-level access to LLMs and avoids the bottleneck of prompt size constraints. Comprehensive experiments across a variety of tabular benchmarks demonstrate that \textsf{FeatLLM} produces robust, high-quality rules, and delivers consistently superior results—achieving average improvements of 10\%—over competing strategies such as TabLLM and STUNT.
\end{abstract}

\section{Introduction}
Recent advances in Large Language Models (LLMs) have showcased their prowess in solving novel tasks with minimal adaptation or fine-tuning~\cite{creswell2022selection,imani2023mathprompter,taylor2022galactica}. By being exposed to vast amounts of textual data during training, LLMs acquire a reservoir of general knowledge, which often serves as an alternative to traditional, domain-specific expertise in diverse settings~\cite{Genomes2021,singhal2023large,wilcox2003role}. This enables substantial progress in automating a wide range of tasks, such as dialogue analysis~\cite{finch2023leveraging} and automated program repair~\cite{fan2023automated}. When given a task description along with a handful of examples, LLMs effectively interpret the underlying context, selectively attending to salient features and instances~\cite{brown2020language,zhao2023survey}. These abilities underline LLMs' promise as sophisticated feature engineering tools capable of data-driven analysis.

The extension of LLMs’ knowledge and reasoning skills has reached the field of tabular learning. For example, prior knowledge such as the association between aging and increased susceptibility to certain diseases can be instrumental for prediction tasks in healthcare. Contemporary research encodes both tasks and features in natural language, serializes the relevant data, and inputs them into LLMs~\cite{dinh2022lift,wang2023anypredict}. Following this, models often adopt either in-context learning techniques~\cite{nam2023semi} or efficient parameter tuning strategies~\cite{dinh2022lift,hegselmann2023tabllm}. The domain expertise captured by LLMs has proven beneficial, particularly in scenarios with scarce training data, consistently outperforming classical tabular learning methods under low-shot conditions~\cite{hegselmann2023tabllm}.

Existing LLM-driven approaches for tabular learning are hampered by several shortcomings. In the case of end-to-end prediction tasks, such methods typically necessitate performing an LLM inference for every individual data point, which leads to significant computational overhead. Additionally, achieving high prediction accuracy often requires fine-tuning the LLM, posing a substantial barrier when training code or tuning APIs are unavailable. This issue is particularly salient as many recent state-of-the-art LLMs restrict users to inference-only API access~\cite{anil2023palm,openai2023gpt4}. Moreover, many LLM-based tabular solutions are ill-equipped to handle extended input prompts~\cite{liu2023lost}; as the number of features grows and the serialized text input lengthens, such methods may become unworkable or their predictive accuracy may degrade. Collectively, these limitations constrain the applicability of LLM-based approaches for practical tabular learning.

To address these limitations, our method reconceptualizes the function of LLMs, harnessing them primarily for feature engineering rather than for end-to-end prediction. We present a novel in-context learning strategy called \textsf{FeatLLM}, which is designed to elucidate the ``criteria" that underlie LLM decision-making. For example, in a disease classification scenario, the LLM can be prompted to articulate explicit rules or conditions linking specific feature patterns to the diagnosis. These extracted rules are then utilized to synthesize new features, which serve as replacements for the original variables in downstream learning tasks. This paradigm offloads predictive complexity from the downstream model to the feature construction stage, thereby simplifying inference and reducing latency since less elaborate models are subsequently required. Capitalizing on LLMs’ in-context learning capabilities, features can be mined without fine-tuning the model, simply by incorporating demonstration examples into the prompt. Additionally, by integrating ensemble learning techniques such as feature bagging~\cite{breiman2001random}, our approach can efficiently deal with issues arising from overly large prompt inputs.

\textsf{FeatLLM} decomposes the process of feature engineering from input prompts into two main components: (1) grasping the nature of the problem and discerning how features relate to target variables, and (2) utilizing both the inferred knowledge from the first step and a limited set of illustrative examples to construct distinguishing rules that separate different classes. This form of guided reasoning encourages LLMs to disregard irrelevant features during rule construction. The rules formulated in this manner are systematically applied to the dataset—interpreting them as programmatic instructions—transforming each data sample into a binary feature that signifies its compliance with each rule. Subsequently, a simple predictive model is trained to infer class probabilities based on these transformed features. To enhance stability, the method integrates an ensemble strategy, generating multiple sets of features alongside bagging techniques. Adjusting the selection of both features and samples within the bagging scheme effectively addresses limitations related to overly large prompts. Notably, \textsf{FeatLLM} requires no additional training and maintains minimal inference overhead, thereby extending the practical usability of LLMs to diverse applied scenarios.

\textsf{FeatLLM} is benchmarked across 13 tabular datasets under low-shot conditions, consistently demonstrating competitive and reliable performance. We demonstrate that LLMs are adept at leveraging pre-existing domain knowledge as well as insights extracted directly from the data, resulting in the generation of informative features. Our proposed framework surpasses current few-shot learning baselines under a range of experimental setups. The corresponding codebase is available through an anonymized GitHub repository at~\texttt{https://github.com/Sungwon-Han/FeatLLM}.

\section{Related Work}

\subsection{Few-Shot Learning with Tabular Data} Progress in few-shot learning has facilitated the derivation of broadly applicable representations from limited annotated data, effectively reducing labeling expenditures~\cite{chen2018closer}. Although this research area began with an emphasis on visual data, the methodologies have since demonstrated remarkable adaptability to multiple modalities, encompassing text, audio, and, more recently, tabular formats~\cite{majumder2022few,nam2023stunt,schick2021exploiting}. Relying on only a handful of labeled samples, however, increases susceptibility to misleading patterns or spurious associations~\cite{bartlett2021deep}. In response, contemporary approaches have integrated supplementary informational sources and incorporated domain-relevant priors to shape model learning through appropriate inductive biases. Examples include leveraging substantial pools of unlabeled data—combined with strategic augmentations—to enable semi-supervised learning~\cite{ucar2021subtab,yoon2020vime}, as well as adopting unsupervised pre-training, independent of the downstream task, to enhance model initialization~\cite{somepalli2021saint}. Unsupervised training strategies often involve objectives like reconstructing partially masked or corrupted data~\cite{arik2021tabnet,majmundar2022met} or utilizing augmented data within contrastive learning frameworks~\cite{bahri2022scarf,chen2020simple}. Nevertheless, the varied structure of tabular datasets complicates the design of generic augmentation techniques, and thus these methods have so far shown limited efficacy under few-shot constraints~\cite{nam2023stunt}.

Recent investigations have explored training algorithms on a diverse assortment of tabular datasets—not restricted to those directly pertinent to the eventual task—and subsequently applying transfer learning on the desired target task~\cite{zhang2023generative,levin2022transfer,zhu2023xtab}. One such example, TransTab~\cite{wang2022transtab}, utilizes transformer-based models to learn semantic linkages among columns, capitalizing not only on the content of table entries but also the semantics of column names. The resulting generalized representations enable zero-shot and few-shot task adaptation to novel tabular tasks. In another line of work, TabPFN~\cite{hollmann2022tabpfn} synthesizes datasets with varying distributions to pre-train a Bayesian neural network; inference on target tasks is then accomplished by simultaneously inputting available training and testing samples, obviating additional training. Incorporation of knowledge from numerous diverse datasets has enabled these approaches to exceed the performance of standard tree-based models, offering considerable advantages in low-data regimes.

\subsection{Language-Interfaced Tabular Learning} An alternative strategy for exploiting prior information is the integration of large language models (LLMs)~\cite{anil2023palm,openai2023gpt4}. LLMs, having been exposed to vast and heterogenous text corpora, exhibit notable generalization to previously unseen problems, affirming their value in many fields~\cite{brown2020language}. Recent efforts have explored harnessing the comprehensive prior knowledge within LLMs for tasks involving tabular data. A common paradigm is to convert tabular inputs into a textual format and inject them, along with relevant feature and task descriptions, into the prompts fed to LLMs~\cite{dinh2022lift,hegselmann2023tabllm,wang2023anypredict}. By supplying such enriched prompts, the models are able to discern connections between tasks and features, thereby facilitating inference. Subsequent advancements have focused on enhancing these models either through parameter-efficient tuning mechanisms such as LoRA~\cite{hu2021lora} or IA3~\cite{liu2022few}, or by adopting in-context learning techniques that incorporate exemplar demonstrations directly in the prompts—without necessitating any parameter updates~\cite{dinh2022lift,hegselmann2023tabllm,nam2023semi,slack2023tablet}.

Although the aforementioned approaches have shown promise in advancing few-shot tabular learning, their dependency on routing all inference operations through the LLM can present practical obstacles, particularly in industrial contexts. The present work seeks to move away from the standard black-box paradigm where every sample is processed end-to-end using an LLM for prediction. Instead, the proposed method leverages the LLM to explicitly deduce the set of rules or conditions defining each class. \textsf{FeatLLM} translates these inferred rules into executable program code that produces novel features. This strategy facilitates predictions without continual LLM involvement, capitalizing on in-context learning to derive rules by utilizing existing domain knowledge as well as the few-shot examples.

\section{Methods}
\textsf{FeatLLM} focuses on uncovering ``rules"—the characteristic conditions defining each answer class—from a limited number of input samples, as illustrated in Figure~\ref{fig:main_model}. Once generated by the LLM, these rules are parsed and utilized to construct new binary features for the data, each indicating whether a particular rule is satisfied by the sample. These newly formed binary features are then combined using a dot product with a set of non-negative, trainable weights to compute class probabilities. Through model training, the relative influence of each rule-based feature on class predictions is learned (refer to Section~\ref{Sec:training}). To enhance robustness, the full process is iteratively applied within a bagging framework, with outcomes aggregated through ensembling. Details regarding the formal problem setup and methodology for each phase are elaborated below.

\vspace*{-0.12in}
\paragraph{Problem formulation.} We examine a tabular dataset $\mathcal{D} = \{(\mathbf{x}^i, \mathbf{y}^i)\}_{i=1}^{N}$ composed of $N$ labeled instances. Each data point $\mathbf{x}^i$ is a vector of dimension $d$, corresponding to $d$ distinct features, and is accompanied by a label $\mathbf{y}^i$ drawn from a finite set of predefined categories, suitable for a classification scenario. The dataset provides feature names in natural language, $F = \{f_j\}_{j=1}^{d}$, and their detailed definitions are available elsewhere. For the purpose of $k$-shot learning evaluations, a subset of $k$ samples, where $k < N$, is randomly selected to train the model.

\subsection{Prompt Design for Extracting Rules}
\label{Sec:prompt_design}
To facilitate rule extraction by the LLM along lines similar to the deductive processes of a domain expert, we design our prompting strategy to emulate human expert reasoning in tabular learning contexts. The overall prompt is structured into three core elements, detailed below and depicted in Fig.~\ref{fig:prompt_for_rule} and Appendix~\ref{sec:example_rule}:

\vspace*{-0.12in}
\paragraph{Basic information description.} This section delivers the critical background necessary for tackling the prediction task (see the orange text in Fig.~\ref{fig:prompt_for_rule}). It encompasses a definition of the target problem with associated label information, comprehensive descriptions of input features, and illustrative examples. The problem is stated as a question (for instance, “Does this patient’s myocardial infarction complications data indicate chronic heart failure? Yes or no?”). Further examples can be found in Appendix~\ref{sec:task_desc}. Feature descriptions specify value types and may also provide formal definitions. For categorical features, sample categories can be listed as illustrations. For the purpose of example demonstration, a handful of labeled training instances are serialized into textual format, making their features and corresponding ground-truth labels explicit. The process of serialization aligns with methodologies established in prior work~\cite{dinh2022lift,hegselmann2023tabllm}:

\begin{align}
&\text{Serialize}(\mathbf{x}^i,\mathbf{y}^i, F) = \nonumber \\ 
&``\ f_1\ \text{is}\ \mathbf{x}^i_1.\ ... \ f_d\  \text{is}\  \mathbf{x}^i_d.\ \ \text{Answer:}\  \mathbf{y}^i\ ”
\end{align}

\vspace*{-0.12in}
\paragraph{Reasoning instruction.} Our objective is to strengthen the LLM's capacity for reasoning by supplying explicit reasoning instructions (illustrated in blue in Fig.~\ref{fig:prompt_for_rule}). These instructions start with an opening line akin to the chain-of-thought strategy~\cite{wei2022chain}. Subsequently, we segment the reasoning undertaken by the LLM into two distinct phases. The initial phase involves prompting the LLM to determine causal links or general tendencies between specific features and the task description, doing so independently of example demonstrations and instead drawing upon its background knowledge and intuition. This approach maximizes the exploitation of pre-existing knowledge relevant to the task. In the following phase, the LLM integrates the insights gathered in the first phase with concrete example demonstrations to synthesize rules pertaining to each class. This bifurcated approach mitigates the risk of the model relying on misleading associations found in non-essential columns and assists in directing attention toward more informative features. To strike a balance and avoid underfitting due to too few rules or unwieldy prompt sizes resulting from an excess of rules, we specify a fixed rule count (i.e., 10)\footnote{Further discussion on rule count is presented in Section~\ref{sec:analysis}.} for extraction per class. When formulating these rules, the LLM is also prompted to consult the feature type information given in the basic description section to ensure rule-feature type consistency.

\vspace*{-0.12in}
\paragraph{Response instruction.} To make the generated rules more readily parsable and practical for downstream use, we instruct the LLM to organize its outputs in a prescribed format (as indicated in yellow in Fig.~\ref{fig:prompt_for_rule}). The prompt specifies formatting conventions both for responses attributed to each answer class (i.e., the response layout) and for the syntactic structure of individual rules (i.e., the rule format). By following these directives, the LLM can correctly adapt the structure based on the corresponding feature types. In the absence of further restrictions, the LLM may utilize logical connectors such as AND or OR to combine several rules. An illustrative result is presented in Appendix~\ref{sec:outcome-rule}.

\subsection{Inferring Class Likelihood via Rules}
\label{Sec:training}

\paragraph{Parsing rules for feature generation.} We employ the rules formulated in the preceding step to generate additional binary features. These newly constructed features are defined for each class, signifying whether a given sample fulfills the criteria for that class’s rules (i.e., the feature takes a value of ‘0’ or ‘1’). These features can subsequently inform predictions regarding a sample's class membership. Nevertheless, because the LLM-derived rules are expressed in natural language, it is necessary to develop an automated parsing method to convert these rules into usable features. Although formatting conventions are imposed during rule generation to facilitate parsing, LLM outputs may still contain syntactic irregularities~\cite{kovavc2023large}. For example, instead of the symbols ``$>$” or ``$<$”, the LLM may output the phrases “greater than” or “less than,” resulting in textual inconsistency and increased parsing complexity.

To mitigate the difficulties arising from parsing such inconsistent output, we refrain from designing intricate parsing programs. Rather, we leverage the LLM’s capacity to comprehend semantics beyond superficial syntactic shifts and its training on code datasets, which enables it to generate simple executable code when provided with clear instructions. Specifically, we prepare prompts that detail the function name, descriptions of inputs and outputs, as well as the set of derived rules, and then submit these prompts to the LLM. Figures~\ref{fig:prompt_for_parsing} and~\ref{sec:example_parsing}-\ref{sec:example_parse_outcome} in the Appendix display example prompts and their resulting outputs. The code generated in this manner is subsequently run using Python’s \texttt{exec()} function to facilitate data transformation.

\vspace*{-0.12in}
\paragraph{Inferring class likelihood.} Once binary features from the rules are constructed per class, a straightforward approach to estimate the likelihood that a sample belongs to a certain class is by tallying the number of satisfied rules for that class (i.e., summing up the corresponding binary features for each class). Nonetheless, since some rules or features are more informative than others, it becomes essential to learn their relative importances from labeled data. We seek to accomplish this by employing a standard machine learning (ML) model for each class’s set of binary features. A linear model without bias, for instance, can be adopted for this purpose. Let $\mathbf{z}_k^i \in \{0, 1\}^R$ denote the binary feature vector for class $k$ from sample $\mathbf{x}^i$, with $R$ indicating the number of rules assigned per class and $c$ denoting the total number of classes. The predicted class probability vector $\mathbf{p}^i$ is calculated as follows:

\begin{align}
\text{logit}_k^i &= \max(\mathbf{w}_k, 0) \cdot \mathbf{z}_k^i, \nonumber \\
\mathbf{p}^i &= \text{Softmax}({[\text{logit}_{1}^i, ..., \text{logit}_{c}^i]}). \label{eq:infer_prob}
\end{align}

In this context, $\mathbf{w}_k$ denote the learnable parameters within the linear model, serving to quantify the significance of each input feature. By constraining these weights to a positive subspace, we guarantee that the features produced by the LLM do not negatively impact the predicted class likelihoods during training. The resulting logit vector is mapped to class probabilities using the softmax function. The cross-entropy loss is subsequently minimized to update the weights $\mathbf{w}_k$, utilizing a limited set of labeled examples during few-shot training.

\vspace*{-0.12in}
\paragraph{Ensembling with bagging.} As a final step, the entire methodology is repeatedly applied to generate a series of inference models whose outputs are then aggregated for final prediction using an ensemble strategy. The ensemble relies on a set of learned weights from $T$ independent trials, represented as $\{\mathbf{w}[t]\}_{t=1}^T$, where the probability for each class $\mathbf{p}^i$ is calculated for each trial's weight $\mathbf{w}[t]$ and subsequently averaged, as expressed in Eq.~\ref{eq:infer_prob}.

To foster variability in the rule generation necessary for a strong ensemble, the temperature parameter during LLM inference is set to a relatively high value. Additionally, the sequence in which few-shot demonstration examples are presented is randomized, recognizing that the order of examples may influence the LLM’s responses~\cite{lu2022fantastically}\footnote{Appendix~\ref{sec:diversity} contains an analysis of rule diversity.}. To further exploit this prediction diversity and thus strengthen model reliability, we employ bagging by randomly selecting a subset of features or data points for each run—a practice that becomes increasingly beneficial with more features or larger training datasets. The ensemble mechanism we propose confers two key benefits. Firstly, should the LLM generate flawed rules in some instances, the ensemble can counterbalance these errors with correct rules derived from other runs, improving the system’s robustness. This phenomenon aligns with the notion of the LLM’s self-consistency~\cite{wang2022self}, where rules frequently produced across several trials are generally more trustworthy. Secondly, our ensembling approach mitigates the limitations imposed by LLM prompt length; as table sizes surpass the permissible prompt length, bagging enables the selection of smaller feature or instance subsets, thus preserving strong predictive performance. While individual rule sets may overlook certain feature interdependencies, aggregating the outputs of multiple weak learners generated across separate trials can offset the gaps arising from partial feature or instance selection in any single model.

\section{Experiments}
We assess the capabilities of \textsf{FeatLLM} across a suite of tabular datasets and perform in-depth analyses to better understand the underlying mechanisms of our framework. Owing to limited space, further details—such as experiments involving alternative LLMs, qualitative case studies, and additional results—are included in the Appendix.

\subsection{Performance Evaluation}
\paragraph{Datasets.}
Our evaluation employs 13 datasets spanning binary and multi-class classification. The first set consists of (1) Adult~\cite{asuncion2007uci}, aimed at classifying individuals by whether their annual income exceeds \$50,000; (2) Bank~\cite{moro2014data}, focused on determining if a client subscribes to a term deposit; (3) Blood~\cite{yeh2009knowledge}, which forecasts donor return behavior; (4) Car~\cite{kadra2021well}, which classifies vehicle quality; (5) Communities~\cite{misc_communities_and_crime_183}, used to estimate crime rates for geographic areas; (6) Credit-g~\cite{kadra2021well}, predicting creditworthiness; (7) Diabetes\footnote{\texttt{kaggle.com/datasets/uciml/pima-indians-diabetes-database}}, for diagnosing diabetes status; (8) Heart\footnote{\texttt{kaggle.com/datasets/fedesoriano/heart-failure-prediction}}, which predicts coronary artery disease; and (9) Myocardial~\cite{misc_myocardial_infarction_complications_579}, evaluating chronic heart failure incidence.

Additionally, our study features several newly released and artificial datasets that are less represented in existing literature and absent from the LLMs' pretraining corpora: (10) Cultivars, designed to predict the expected grain yield of soybean cultivars\footnote{\texttt{archive.ics.uci.edu/dataset/913}}; (11) NHANES, where the aim is to infer the age group of individuals from their records\footnote{\texttt{archive.ics.uci.edu/dataset/887}}; (12) Sequence-type, which involves categorizing numerical sequences as arithmetic, geometric, Fibonacci, or Collatz; and (13) Solution-mix, a binary classification task to determine whether a mixture's concentration from four components surpasses 0.5. Cultivars and NHANES were added to the UCI Machine Learning Repository after September 2023, ensuring exclusion from the pretraining dataset of the LLM API (gpt-3.5-turbo-0613) used in our experiments. The final two datasets are synthetically generated, guaranteeing they were not present during the API’s training phase.

A summary of dataset sizes and complexities is provided in Table~\ref{Tab:basic-desc}. Notably, each dataset is well-documented, offering descriptive and interpretable attribute names. Datasets such as `Communities' and `Myocardial' are high-dimensional, each containing in excess of 100 features, presenting potential hurdles in light of LLMs’ finite context windows.

\vspace*{-0.12in}
\paragraph{Baselines.}
Our approach is compared against nine established baselines. The first set comprises standard tabular prediction methods: (1) Logistic Regression (LogReg), (2) XGBoost~\cite{chen2016xgboost}, and (3) RandomForest~\cite{ho1995random}. Two additional baselines utilize unlabeled data: (4) SCARF~\cite{bahri2022scarf} and (5) STUNT~\cite{nam2023stunt}. It is important to highlight that amassing substantial volumes of unlabeled data can be challenging in applied scenarios. As a more recent approach, (6) TabPFN~\cite{hollmann2022tabpfn} leverages extensive pretraining on synthetic data generated with diverse statistical distributions. The final group consists of LLM-driven strategies: (7) In-context learning (In-context)~\cite{wei2022emergent}, (8) TABLET~\cite{slack2023tablet}, and (9) TabLLM~\cite{hegselmann2023tabllm}. In-context learning integrates a small set of labeled examples within the input prompt without additional parameter updates. TABLET enriches LLM prompts by appending external signals, such as rule-based or prototype-based hints, via a supporting classifier to refine decision quality. TabLLM introduces efficient parameter fine-tuning through IA3~\cite{liu2022few} applied to few-shot data to enhance LLM adaptation.

\vspace*{-0.12in}
\paragraph{Implementation details.}
\textsf{FeatLLM} adopts GPT-3.5 as its foundational LLM, but its architecture remains flexible regarding the underlying language model (refer to Appendix~\ref{sec:backbones} for evaluations with alternative LLM backbones). During inference, we configure the LLM with a temperature of 0.5 and retain the API's default top-p value of 1. The ensemble count and the number of extracted rules are set to 20 and 10, respectively. Analyses of hyper-parameter effects can be found in Figure~\ref{fig:hyperparameter2} and Appendix~\ref{sec:hyper-parameter}. For the linear classifier, training proceeds for 200 epochs utilizing the Adam optimizer with a learning rate fixed at 0.01. We apply $k$-fold cross-validation to determine the optimal epoch for model selection. When reproducing baselines, GPT-3.5 supports in-context learning, while TabLLM uses T0~\cite{sanh2022multitask}, consistent with its original configuration. It is important to highlight that TabLLM limits LLM usage to those with open checkpoint availability; thus, GPT-3.5 is excluded. For classical machine learning approaches—including LogReg, XGBoost, and RandomForest—hyperparameters were chosen via a Grid-search procedure integrated with $k$-fold cross-validation, with $k$ set to 2 or 4 to ensure each training fold contains at least one sample per class. Further implementation specifics are provided in Appendix~\ref{sec:baselines}.

\vspace*{-0.12in}
\paragraph{Results.}
Tables~\ref{tab:main1} and \ref{tab:main2} provide comparative performance metrics across diverse datasets. All experiments are run three times with varying random splits for training and testing, reporting the mean and standard deviation of the resulting AUC values. \textsf{FeatLLM} persistently attains either the highest or second-highest ranking among baseline methods. Remarkably, our method reaches similar levels of accuracy to established tabular models requiring 32 shots, while using only 4 shots (refer to Fig.~\ref{fig:compares}). It is worth mentioning that although STUNT and TabPFN achieved marked improvements on selected datasets, their aggregate advantage over traditional machine learning algorithms was limited. Additionally, LLM-based methods such as in-context learning, TABLET, or TabLLM, were unable to perform inference on tabular datasets with a high dimensionality. Leveraging bagging and ensembling, our framework remains effective and achieves strong results even when handling datasets with many features. It should be emphasized that our bagging and ensembling methodology could be integrated into other LLM-based techniques; however, this would double the inference cost per sample, significantly raising computational requirements and rendering such a modification impractical.

\subsection{Analysis \& Discussion}
\label{sec:analysis}
\paragraph{Ablation study.} To evaluate the significance of each major component, we conduct a series of ablation studies by excluding: (1) \textsf{-Tuning}: the linear model's weight tuning phase, replacing it with a simple sum of binary features for logit calculation; (2) \textsf{-Ensemble}: the ensembling mechanism; (3) \textsf{-Description}: the feature description; and (4) \textsf{-Reasoning}: the Step-1 process in the reasoning instruction used during rule generation.

Table~\ref{tab:ablation} displays the average variation in AUC following each ablation across all datasets. Removal or modification of any component consistently leads to inferior results. Of particular note, the contribution of weight tuning grows as the number of shots increases, suggesting that with ample data, more accurate rule weighting has greater impact. Conversely, when fewer shots are available, incorporating feature descriptions and reasoning instructions delivers substantial benefits, indicating that effective exploitation of LLMs' prior knowledge becomes especially important in low-data settings. Across all scenarios, ensembling reliably enhances performance.

\vspace*{-0.12in}
\paragraph{Training \& inference time comparison.} We evaluate and contrast the training and inference times across various approaches. All experiments utilize the Adult dataset, with each method trained on a subset comprising 16 samples. A single A100 GPU is used for all models, with the exception of TabLLM, which relies on four A100 GPUs for distributed model parallelism. Table~\ref{tab:cost} presents the measured runtimes. Of particular interest, \textsf{FeatLLM} demonstrates inference efficiency that is on par with standard baselines. Conversely, LLM-based models like In-context learning, TABLET, and TabLLM incur substantially longer inference durations. With respect to training overhead, \textsf{FeatLLM} exhibits similar time requirements to its LLM-based counterparts. Crucially, it only issues 30 API calls, which is both cost-efficient and amenable to parallel execution, minimizing budgetary concerns.

\vspace*{-0.12in}
\paragraph{Quality of features generated by \textsf{FeatLLM}.}
We perform a straightforward study to evaluate the informativeness of the rule-derived features produced by \textsf{FeatLLM} for real-world problems. For this purpose, we calculate the mean AUROC between each newly created feature and the ground-truth labels, subsequently determining the proportion of features achieving an AUROC exceeding 0.5 within each dataset. 
A summary of these findings is displayed in Table~\ref{tab:quality}. The evidence shows that the bulk of constructed rules bear strong relevance to the prediction target and yield meaningful contributions for task resolution (see the “Before tuning” column). Furthermore, during linear model optimization, features with marginal importance (i.e., those assigned weights below a certain cutoff) are systematically pruned (refer to the “After tuning” column). The cutoff threshold for weights is fixed at 0.1, informed by the aggregate weight distribution.

\vspace*{-0.12in}
\paragraph{Effect of adding spurious correlations.} To evaluate the robustness of different approaches in eliminating irrelevant spurious correlations under low-shot conditions, we designed a specific analysis. In this study, we employed the Heart dataset and artificially augmented it by appending randomly selected columns from the Adult dataset that do not pertain to the prediction task at hand. We assessed model performance as the number of unrelated features increased incrementally. To ensure a consistent basis for evaluation, we set the labeled data to 8 shots for all models, intentionally omitting any unlabeled data and therefore not including approaches such as SCARF and STUNT in the comparison.

Figure~\ref{fig:spurious} presents the impact of added spurious correlations on model accuracy. Among the evaluated models, \textsf{FeatLLM} exhibited the smallest decline in performance in the presence of such noisy features. This outcome may be attributed to the design of our approach, which leverages prior domain knowledge to link tasks with relevant features during the rule extraction process, discouraging the creation of rules dependent on irrelevant columns and promoting model robustness. Notably, the proportion of rules formed around these noisy columns remained limited to just 1-2\% of the total. However, our analysis further reveals that \textsf{FeatLLM} is not immune to learning misleading associations: when features from data sources closely related to the target are randomly incorporated, the model may also encode spurious correlations and misconstrue causal dependencies between the target and features (refer to Appendix~\ref{sec:spurious-appendix}).

\vspace*{-0.12in}
\paragraph{Hyper-parameter analysis}
\textsf{FeatLLM} introduces two principal hyper-parameters: the ensemble count and the number of rules generated per class in each LLM query. We systematically evaluated how variations in these hyper-parameters influence model outcomes. The experiments indicated that increasing the ensemble size generally boosts performance, but the improvement plateaus when the ensemble number exceeds approximately 20 (refer to Fig.~\ref{fig:appendix-hyperparameter1} in Appendix). Regarding the rule count, we observed no statistically significant variations in results; however, setting this parameter to 10 offered the optimal balance between prompt size and predictive effectiveness (see Fig.~\ref{fig:hyperparameter2}). We hypothesize that extracting more rules increases the dimensionality of the input, introducing richer information but at the cost of heightened risk for overfitting, especially within the context of low-shot scenarios.

\section{Conclusion}
We introduce \textsf{FeatLLM}, a method that raises the bar for LLM-driven few-shot tabular learning. Unlike approaches that restrict LLM usage to per-sample predictions, \textsf{FeatLLM} exploits the generative capacity of LLMs to develop decision rules that emulate feature engineering. By leveraging a rule-based approach and aggregating predictions via bagging, \textsf{FeatLLM} is computationally efficient at inference, relies solely on API access without necessitating model retraining, and operates without strict limitations on the number of data features. The framework yields robust predictive performance, making it a compelling and resource-efficient choice for real-world applications of LLMs in tabular data contexts.

The design of \textsf{FeatLLM} is specifically tailored to settings with few available labeled samples. Looking ahead, we plan to enhance its applicability to larger datasets by enabling the automated generation of novel features, incorporating feature construction methods that go beyond rule induction, and embedding interpretability mechanisms, thereby broadening its utility across diverse machine learning problems.

\section{Broader Impact}
With \textsf{FeatLLM}, we propose a framework that capitalizes on the vast prior knowledge embedded within LLMs to extend the capabilities of conventional supervised tabular learning, specifically by boosting data efficiency. By integrating prior knowledge, our approach provides a solution to overfitting challenges often encountered when labeled examples are scarce. This is especially advantageous for stakeholders in domains such as finance and healthcare, where annotation processes are resource-intensive or complicated, and where LLM pretraining with domain-relevant corpora can impart considerable benefits. As a critical direction for future exploration and widespread deployment, we recognize the need to scrutinize the influence of inherent societal biases or inaccuracies that may be present in the LLMs’ pretraining data, as these could profoundly shape model predictions—potentially superseding the effect of ground-truth labeled data.

\end{document}